\subsection{Activation Pattern Table Design}

After the OTFS-IM block parameters have been defined, the next design step is to determine how the index bits are converted into active delay-Doppler positions. This section describes the activation-pattern table used by the transmitter and receiver. The purpose is not only to list possible patterns, but also to explain why the pattern table has a direct influence on index detection reliability.

\subsubsection{Constant-Weight Pattern Table}

In each OTFS-IM block, exactly \(k\) out of \(n\) delay-Doppler positions are active. Therefore, every valid activation pattern has the same Hamming weight \(k\). The complete set of possible activation patterns is generated by choosing \(k\) positions from \(n\) available positions:
\begin{equation}
    \mathcal{A}_{\mathrm{all}}
    =
    \left\{
    \mathbf{a} \in \{0,1\}^{n}
    :
    \sum_{i=1}^{n} a_i = k
    \right\}.
    \label{eq:all_constant_weight_patterns}
\end{equation}
The number of candidate patterns is \(\binom{n}{k}\). Each candidate can be stored either as a set of active positions or as a binary vector. For example, when \(n=4\) and \(k=3\), a pattern row \([1,2,4]\) means that the first, second, and fourth positions in the local IM block are active, while the third position is inactive.

For distance evaluation, the position-index representation is converted into a binary vector representation. If the active positions are \([1,2,4]\), the corresponding binary pattern is
\begin{equation}
    \mathbf{a} = [1,1,0,1].
\end{equation}
This binary representation is useful because the difference between two activation patterns can be measured directly by their Hamming distance.

Although \(\binom{n}{k}\) patterns may be available, not all of them are necessarily used. The number of index bits per block is
\begin{equation}
    b_1 =
    \left\lfloor
    \log_2 \binom{n}{k}
    \right\rfloor ,
\end{equation}
so only \(2^{b_1}\) patterns can be assigned to binary input labels. The selected pattern table is denoted by
\begin{equation}
    \mathcal{P}
    =
    \left\{
    \mathbf{p}_0, \mathbf{p}_1, \ldots, \mathbf{p}_{2^{b_1}-1}
    \right\},
    \label{eq:selected_pattern_table}
\end{equation}
This table is shared by the transmitter and receiver. The transmitter uses it to decide which delay-Doppler positions are active, while the receiver uses the same table as the set of candidate patterns in the block-wise detection process.

\subsubsection{Index-to-Pattern Mapping}

In each IM block, the first \(b_1\) bits are interpreted as index bits. These bits are converted into a decimal index and then mapped to one row of the selected table \(\mathcal{P}\). The binary index is first obtained as
\begin{equation}
    q =
    \mathrm{bin2dec}
    \left(
    \mathbf{b}_{\mathrm{idx}}
    \right),
    \qquad
    q \in \{0,1,\ldots,2^{b_1}-1\}.
    \label{eq:index_binary_decimal}
\end{equation}
The simulation then applies Gray indexing before selecting the pattern row:
\begin{equation}
    m = q \oplus \left\lfloor \frac{q}{2} \right\rfloor ,
    \label{eq:gray_index_mapping}
\end{equation}
where \(\oplus\) denotes bitwise exclusive-or. The active positions for the current block are then read from
\begin{equation}
    \mathbf{p}_m = \mathcal{P}[m].
\end{equation}

For the \(i\)-th IM block, the local active positions \(\mathbf{p}_m\) are shifted to the corresponding positions in the vectorized delay-Doppler frame. The QAM symbols are inserted only at these active locations, while the other positions in the same block remain zero. Therefore, the transmitted block carries information through both the selected pattern and the QAM symbols placed on the active positions.

At the receiver, the same pattern table is used in the reverse direction. The detector evaluates every candidate row of \(\mathcal{P}\) and selects the pattern that gives the largest likelihood score. After the most likely Gray index is found, it is converted back to the original binary index bits. This ensures that the transmitter and receiver use a consistent mapping between index bits and activation patterns.

The use of a fixed pattern table is important for reproducibility. If the transmitter and receiver used different pattern ordering, the active positions might still be detected correctly but the recovered index bits would be wrong. Therefore, the pattern table must be treated as part of the system design rather than as a temporary implementation detail.

\subsubsection{Pattern-Distance Interpretation and Selection Need}

The reliability of index detection depends strongly on how distinguishable the selected activation patterns are. Two patterns that differ in only a small number of positions are more likely to be confused after channel distortion, noise, and imperfect probability estimation. This effect is measured by the Hamming distance between two binary patterns:
\begin{equation}
    d_{\mathrm{H}}(\mathbf{p}_i,\mathbf{p}_j)
    =
    \sum_{\ell=1}^{n}
    \left|
    p_i(\ell) - p_j(\ell)
    \right| .
    \label{eq:pattern_hamming_distance}
\end{equation}
For constant-weight patterns, the Hamming distance counts how many local block positions change their active or inactive status between two candidate patterns. A larger Hamming distance generally means that the receiver must make more activation-position mistakes before one pattern is confused with another.

The minimum pairwise Hamming distance of a selected pattern table is defined as
\begin{equation}
    d_{\min}
    =
    \min_{i \ne j}
    d_{\mathrm{H}}(\mathbf{p}_i,\mathbf{p}_j).
    \label{eq:pattern_min_distance}
\end{equation}
This metric is useful because the worst-separated pair of patterns often dominates the pattern error behavior. If two patterns in the table are too similar, a small detection error in one delay-Doppler position can cause an incorrect index decision.

However, maximizing only the number of possible patterns is not always the best design choice. Since only \(2^{b_1}\) patterns are used, the system has an opportunity to select a subset with better distance properties. This is the motivation for the OHD-based pattern selection methods discussed later in Section 4.6. In the present section, the important point is that \(\mathcal{P}\) is not merely a lookup table for implementation convenience. It determines the geometric separation among index labels and therefore affects index BER, pattern error rate, and total OTFS-IM BER.

In the simulation results of Chapter 5, this effect is evaluated by comparing different pattern-selection methods under the same \((n,k)\), modulation order, and channel setting. The total BER alone is not sufficient to explain the role of the activation table, so index BER and pattern error rate are also reported. These additional metrics reveal whether performance changes come from symbol detection errors or from incorrect activation-pattern decisions.
